\subsection{Restrições dos Atuadores}

Além dos limites cinemáticos e geométricos das juntas e dos elos do manipulador robótico (MR), ele está sujeito às restrições dinâmicas implementadas nos atuadores, que são expressas em termos de torque máximo nas juntas revolutas e força máxima na junta prismática.  
As restrições são expressas conforme:
\begin{equation}
\label{eq:limite_torque}
|\tau_j| \leq \tau_{j,max} \qquad \text{(juntas revolutas)},
\end{equation}

\begin{equation}
\label{eq:limite_forca}
|f_5| \leq f_{d5,max} \qquad \text{(junta prismática)}.
\end{equation}

% \begin{equation}
% \label{eq:limit_forca}
% |f_5(t)| \leq f_{5,\max}, \qquad
% |\tau_i(t)| \leq \tau_{i,\max}, \;\; i=1,\dots,4.
% \end{equation}

Esses valores máximos foram definidos de forma arbitrária, utilizando como base as simulações realizadas.
% , e os valores de $\tau_{j,\text{max}}$ e $f_{d5,max}$ serão apresentados posteriormente com base na análise dinâmica completa e na sintonia do controlador implementado no ambiente de simulações \emph{MATLAB/Simulink}.  

